\section{Introduction}
\label{sec:intro}

Cross-country cryptocurrency price gaps are often treated as direct evidence of segmented digital-asset markets, capital controls, or unexploited arbitrage \citep{pieters2017financial,choi2018kimchi,makarov2020trading,borri2020regulation,grafvonluckner2023crypto,hautsch2024building}. In emerging markets with managed-float currencies, however, domestic spot exchanges settle in local fiat while global reference venues trade in USD-pegged stablecoins such as Tether (USDT). Measured against Bank Indonesia's official Jakarta Interbank Spot Dollar Rate (JISDOR), Indonesian bitcoin on Indodax traded at an average premium of 30 basis points between January 2021 and November 2026. This paper asks whether that premium reflects mispricing of bitcoin itself or a domestic rupiah-to-stablecoin conversion spread.

Indonesia is a clean setting for separating crypto-market pricing from local currency and stablecoin frictions. Licensed domestic spot exchanges (Pedagang Fisik Aset Kripto), led by Indodax, quote directly in Indonesian rupiah (IDR), while global price discovery concentrates in USDT pairs on offshore venues such as Binance. Between January 2021 and November 2026, Indonesia also saw major institutional changes. Bank Indonesia shifted JISDOR's publication schedule in April 2021 and again in April 2026 \citep{bi_sp233021,bi_padg072026}. The Ministry of Finance introduced crypto VAT and final income taxes under PMK No.\ 68/PMK.03/2022 and replaced them in August 2025 under PMK No.\ 50 Tahun 2025 \citep{pmk68_2022,pmk50_2025}. Supervision moved from BAPPEBTI to the Financial Services Authority (OJK) in January 2025 alongside the launch and later fee cuts of the domestic Crypto Exchange (CFX) and clearing house (KKI). Earlier empirical studies claiming persistent Indonesian crypto segmentation \citep{maulana2026segmentasi} lacked fixing-synchronized intraday data and regime-accurate execution costs to separate these channels.

We build a fixing-synchronized panel covering January 4, 2021 to November 12, 2026 ($n = 1,414$ JISDOR business days across BTC, ETH, and XRP), paired with continuous five-minute data ($N = 616,683$ bars) and 2026 one-minute data ($N = 455,040$ bars). We match Indodax rupiah books against Binance USDT spot klines, venue-specific Indodax \texttt{USDT/IDR} candles, and Coinbase \texttt{USDT-USD} for the external USD/USDT peg wedge. Pre-November 2025 evidence rests on Indodax alone; for the final 349 days we cross-check Indodax directly against Binance's native IDR pairs.

Our results extend the fiat-leg interpretation of cross-country crypto spreads in \citet{makarov2020trading} to a managed-float emerging market with an active domestic stablecoin book. First, once domestic bitcoin prices are converted via local \texttt{USDT/IDR}, the pure crypto leg averages $-0.83\text{ bps}$ (statistically indistinguishable from zero), and $98.44\%$ of business-day deviations sit inside a reconstructed schedule of Indonesian taxes, CFX/KKI levies, and exchange taker fees ($60.0$ to $81.0\text{ bps}$ across regimes). Second, by accounting identity, total BTC basis variance relative to JISDOR splits into $83.10\%$ from the domestic \texttt{USDT/IDR} minus JISDOR spread, $11.53\%$ from the pure crypto leg, and $5.37\%$ from their covariance. Because this split is an identity, the shares describe variance rather than a causal channel; they are also not an artifact of pairing an all-day FX average with a point snapshot, as volume-averaging \texttt{USDT/IDR} over JISDOR's full intraday monitoring window leaves the stablecoin--FX share at $82.39\%$.

Third, after Holm multiple-testing adjustment across ten institutional events and both basis legs, only the August 2025 PMK 50/2025 tax overhaul shifts the pure crypto leg ($+21.23\text{ bps}$, $p_{\text{Holm}} = 0.011$), consistent with its shift of statutory tax incidence from buyers to sellers. Finally, daily error-correction estimates suffer from near-collinear cross-market shocks, making daily Hasbrouck information shares unidentifiable. Moving to a five-minute VECM converted with synchronous Indodax \texttt{USDT/IDR} lowers residual correlation to $0.671$ and shows that local \texttt{BTC/IDR} does most of the short-run error correction (half-life $11.04$ five-minute bars; Gonzalo--Granger global weight $83.6\%$).

Section~\ref{sec:institutional} documents Indonesian FX fixing mechanics, tax law changes, CFX fees, and platform commission regimes. Section~\ref{sec:data} details data sources, coverage limits, and synchronization rules. Section~\ref{sec:method} sets out the three-leg basis identity, cointegration tests, VECM, price discovery, and threshold models. Section~\ref{sec:results} reports the empirical findings. Section~\ref{sec:robustness} presents synchronization sensitivity, JISDOR-window averaging checks, and falsification tests. Section~\ref{sec:discussion} covers limitations and policy implications, and Section~\ref{sec:conclusion} concludes.
